\section{Estudio de caso: control de robots}

\subsection{Offline policy in robotics}

En las tareas de robotics, el costo de recopilar datos es alto y el offline RL puede reducir considerablemente los real-world trials. Práctica típica:

\begin{itemize}
    \item Entrenar la policy en el simulator y hacer replay de las traj clave en el physical robot.
    \item Usar el reward multiplier del offline dataset para enfatizar las acciones relacionadas con la seguridad (como low velocity).
    \item Registrar las sensor readings de each deployment run para poder hacer un counterfactual audit de los errores.
\end{itemize}

\begin{table}[H]
\centering
\begin{tabular}{p{0.25\textwidth} p{0.2\textwidth} p{0.2\textwidth} p{0.2\textwidth}}
\toprule
Task & Data source & Safety guardrail & Offline method \\
\midrule
Arm manipulation & Robot logs + human demos & Collision penalties & CQL \\
Drone navigation & Simulation + recorded flights & Wind model alert & IQL + replay \\
\bottomrule
\end{tabular}
\caption{Configuración de offline RL en escenarios de robots}
\end{table}

\subsection{Resumen de la sección}

El estudio de caso de robots destaca que son indispensables: la diversidad de data source, el reward carefully shaped y el monitoring & replay.

